\documentclass[12pt,a4paper]{article}
\usepackage[T2A]{fontenc}
\usepackage[utf8]{inputenc}
\usepackage[russian]{babel}
\usepackage{amsmath,amssymb,amsfonts}
\usepackage{geometry}
\usepackage{titlesec}
\usepackage{enumitem}
\usepackage{booktabs}
\usepackage{xcolor}
\usepackage{hyperref}
\usepackage{mathtools}
\usepackage{tikz}
\usepackage{pgfplots}
\pgfplotsset{compat=1.18}
\usetikzlibrary{calc, arrows.meta}
\usetikzlibrary{patterns, fillbetween, intersections}
\usepackage{fancyhdr}
\geometry{left=2.5cm,right=2.5cm,top=2.5cm,bottom=2.5cm}
\pagestyle{fancy}
\fancyhf{}
\fancyhead[L]{Машинное обучение 1 --- Контрольная 2024 --- Лашкевич Злата БЭАД245}
\fancyfoot[R]{\thepage}
\renewcommand{\headrulewidth}{0.4pt}
\renewcommand{\footrulewidth}{0pt}
\setlength{\headheight}{14.5pt}
\newcommand{\defeq}{\overset{\mathrm{def}}{=\joinrel=}}

\begin{document}
\selectlanguage{russian}
\section*{Задача 1 (2.5 балла). Функции потерь и SVM}
Для начала напомним, как выглядит кусочно-постоянная функция потерь из SVM:
\[
L(y,z) = \max(0, 1-yz)
\]
Ответьте на вопросы о функциях потерь и линейных моделях:

\begin{enumerate}
\item В функции потерь для SVM отсечка в точке $yz=1$ беспокоит нас. Мы решили поэкспериментировать и воспользоваться функцией потерь $\tilde{L} = \max(0, -1 - yz)$. Что получится, если мы обучим модель на такую функцию?

\item Мы продолжаем эксперименты и хотим сделать два <<излома>>: $\hat{L} = \min(\max(0, 1-yz), 10)$. Что будет при обучении на эту функцию потерь?

\item Рассмотрим обобщённую задачу метода опорных векторов:
\begin{equation}
\begin{cases}
\|w\|^p + C\displaystyle\sum_{i=1}^{\ell}\xi_i^q \to \min_{w,b,\xi} \\[6pt]
y_i(\langle w, x_i\rangle + b) \geqslant 1 - \xi_i, \quad i=1,\ldots,\ell, \\[4pt]
\xi_i \geqslant 0, \quad i=1,\ldots,\ell.
\end{cases}
\tag{0.1}
\end{equation}
Считайте, что $p>0$, $q>0$. Ответьте на вопросы:
\begin{itemize}
\item Запишите эту задачу оптимизации в безусловном виде.
\item Как гиперпараметры $p$ и $q$ влияют на задачу? Опишите, как будет меняться решение при уменьшении и увеличении каждого из этих гиперпараметров.
\end{itemize}
\end{enumerate}

\subsection*{1.}

Введём обозначение отступа: $M = yz$. Стандартный Hinge loss записывается через отступ как $L(M) = \max(0, 1 - M)$ и равен нулю при $M \geqslant 1$, то есть модель перестаёт штрафовать объекты, классифицированные верно с отступом не менее 1.\\[1mm]
Новая функция потерь:
\[
\tilde{L}(M) = \max(0, -1 - M)
\]

Эта функция равна нулю при $M \geqslant -1$. Значит, модель не штрафует объекты с отступом $M \in [-1, +\infty)$.

\textbf{Как следствие} все объекты, классифицированные верно ($M > 0$), а также объекты с небольшими ошибками (отрицательный отступ $M \in [-1, 0)$) не вносят никакого вклада в функционал ошибки. Модель начинает штрафовать только объекты с $M < -1$, то есть те, которые классифицированы неверно и при этом находятся далеко за разделяющей гиперплоскостью.

Хорошая функция потерь должна штрафовать за неправильные ответы и за неуверенную классификацию. Здесь же модель допускает ошибки с отступом до $-1$ без какого-либо штрафа, что означает, что разделяющая гиперплоскость может проходить так, что часть объектов будет классифицирована неверно, но модель не будет стремиться это исправить.

\textbf{Вывод:} При обучении на $\tilde{L}$ модель будет существенно хуже разделять классы по сравнению со стандартным SVM, поскольку допускает грубые ошибки без штрафа.

\subsection*{2.}

Запишем через отступ $M = yz$:
\[
\hat{L}(M) = \min\bigl(\max(0, 1-M),\; 10\bigr)
\]

Поведение по трём областям:
\begin{center}
\begin{tabular}{cll}
\toprule
\textbf{Область} & \textbf{Значение $\hat{L}$} & \textbf{Производная} \\
\midrule
$M \geqslant 1$        & $0$       & $0$ \\
$-9 \leqslant M < 1$   & $1-M$     & $-1$ \\
$M < -9$               & $10$      & $0$ \\
\bottomrule
\end{tabular}
\end{center}

При $M \leqslant -9$ (грубо неверно классифицированные объекты) штраф перестаёт расти и фиксируется на уровне 10. Градиент в этой области равен нулю.

Функция потерь считается устойчивой к выбросам, если её производная ограничена при $M \to -\infty$. У стандартного Hinge loss производная равна $-1$ при $M < 1$ --- она ограничена, но штраф растёт неограниченно. У $\hat{L}$ штраф ограничен числом 10.

\textbf{В следствие этого} модель становится \textit{устойчивой к сильным выбросам}. Объекты с очень большим отрицательным отступом ($M < -9$) не влияют на веса при обучении, из-за нулевого градиента. Обычный SVM пытался бы существенно изменить веса, чтобы уменьшить огромный штраф на таких объектах. Ограниченный SVM просто <<заплатит>> фиксированный штраф 10 и проигнорирует выброс, сосредоточившись на основной массе данных.

\subsection*{3.}

\subsubsection*{Безусловный вид}

Из ограничений задачи (0.1):
\[
\xi_i \geqslant 1 - y_i(\langle w, x_i\rangle + b), \qquad \xi_i \geqslant 0
\]
Поскольку мы минимизируем $\sum \xi_i^q$ (а $q > 0$), то оптимальное значение:
\[
\xi_i = \max\bigl(0,\; 1 - y_i(\langle w, x_i\rangle + b)\bigr)
\]
Подставляем в целевую функцию:
\[
\|w\|^p + C\sum_{i=1}^{\ell}\Bigl(\max\bigl(0,\;1 - y_i(\langle w, x_i\rangle + b)\bigr)\Bigr)^q
\;\to\;
\min_{w,\,b}
\]

\subsubsection*{Влияние $p$}

Параметр $p$ определяет тип регуляризации:

\begin{itemize}
\item $p = 2$: стандартная $L_2$-регуляризация. В SVM это эквивалентно максимизации ширины разделяющей полосы $2/\|w\|$. Веса стремятся быть небольшими, но ненулевыми

\item $p = 1$: $L_1$-регуляризация. Приводит к разреженности весов --- многие компоненты вектора $w$ обнуляются. Это реализует встроенный отбор признаков

\item При $p \to 0$: штрафуется количество ненулевых весов (псевдонорма $L_0$). Модель стремится использовать минимальное число признаков
\end{itemize}
При увеличении $p$ регуляризатор становится менее чувствительным к малым и средним весам, модель допускает больше таких весов, что может привести к переобучению. Большие веса штрафуются ещё сильнее, поэтому большие $p$ приводят к решениям с равномерно небольшими, но ненулевыми весами.  А при уменьшении $p$ (особенно $p < 1$) регуляризатор всё сильнее поощряет разреженность, но задача оптимизации становится невыпуклой

\subsubsection*{Влияние $q$}

Параметр $q$ определяет степень штрафа за нарушение отступа:

\begin{itemize}
\item $q = 1$: стандартный Hinge loss. Линейный штраф за нарушение отступа. Производная ограничена, функция устойчива к выбросам

\item $q = 2$: квадратичный Hinge loss. Сильнее штрафует за большие ошибки: если отступ нарушен на 2, то при $q=1$ штраф равен 2, а при $q=2$ штраф равен 4. Менее устойчив к выбросам

\item $q < 1$ (например, $q = 0.5$): штраф растёт медленнее линейного. Производная $q\xi^{q-1}$ стремится к бесконечности при $\xi \to 0$, но при больших $\xi$ штраф растёт медленно. Это повышает устойчивость к объектам с большими нарушениями отступа
\end{itemize}
При увеличении $q$ модель всё сильнее подгоняется под выбросы, поскольку квадрат и более высокие степени ошибки доминируют в функционале (аналогия с MSE vs MAE). При уменьшении $q$ модель становится более робастной, аналогично тому, как MAE устойчивее MSE к выбросам, но задача теряет выпуклость


\section*{Задача 2 (2.5 балла)}

Ответьте на вопросы по решающим деревьям и композициям:

\begin{enumerate}
\item На лекциях обсуждалось, что решающие деревья можно использовать для генерации нелинейных признаков: берём дерево и используем номер листа, куда попал объект, как значение нового категориального признака. При этом можно воспользоваться деревом, обученным привычным нам способом на всей обучающей выборке, а можно взять случайное дерево (где все предикаты выбирались случайно, без оптимизации какого-либо критерия). Удивительно, но последний способ используют чаще при генерации признаков. Чем это можно объяснить? А что нужно сделать с обучаемыми деревьями, чтобы их тоже было удобно использовать для генерации признаков?

\item Иногда в Random Forest используют дополнительную рандомизацию: при построении вершины выбирается (как обычно) случайное подмножество признаков, для каждого признака оттуда генерируется случайный набор порогов, и лучший предикат ищется только среди этих признаков и порогов. Как это изменение влияет на смещение и разброс композиции? Обоснуйте ваш ответ (здесь предполагаются рассуждения, а не строгий формальный вывод).

\item Представьте, что мы решили обучать градиентный бустинг так:
\[
\frac{1}{\ell}\sum_{i=1}^{\ell}\bigl(b_N(x_i) - s_i[|s_i| \geqslant 0.1]\bigr)^2 \to \min_{b_N};\quad
s_i = -\left.\frac{\partial L(y_i, z)}{\partial z}\right|_{z=a_{N-1}(x_i)}
\]
Для конкретики будем считать, что используется логистическая функция потерь $L(y,z) = \log(1+\exp(-yz))$. Как эта модификация повлияет на обучение для первых базовых моделей в композиции? А для более поздних базовых моделей?
\end{enumerate}

\subsection*{1.}

Номер листа дерева $T_n(x)$ является ценным категориальным признаком, который задаёт нелинейное разбиение пространства и несёт информацию о сходстве объектов.\\[1mm]
\textbf{Почему случайные деревья используют чаще:}

Случайные деревья строятся очень быстро, так как не нужно оптимизировать критерий информативности и при этом генерируют \textit{разнообразные} слабо коррелированные признаки за счёт случайности структуры, так каждое случайное дерево даёт своё уникальное разбиение пространства.

Обучаемые деревья, наоборот, \textit{жадные}: при обучении на одной и той же выборке они будут строить очень похожие разбиения, оптимизируя одни и те же критерии. Дисперсия ансамбля:
\[
\text{Var}(a) = \rho\,\sigma^2 + \frac{1-\rho}{N}\sigma^2
\]
корреляция $\rho$ между обучаемыми деревьями будет высокой, а разнообразие признаков --- низким.

Более того, обучаемые деревья, построенные на той же выборке, на которой потом обучается финальная модель, содержат информацию о целевой переменной. Это приводит к <<утечке>> целевой переменной в значения признаков как со стекингом, что может вызвать переобучение.\\[1mm]
\textbf{Что нужно сделать с обучаемыми деревьями:}

Чтобы обучаемые деревья давали разнообразные и некоррелированные признаки, к ним нужно применить рандомизацию:
\begin{itemize}
\item Обучать каждое дерево на бутстрапированной подвыборке как в бэггинге
\item Использовать случайный выбор признаков при каждом разбиении как в Random Forest
\item Обучать деревья на одной части выборки, а использовать признаки --- на другой (аналогия с кросс-валидацией при стекинге)
\end{itemize}

Это уменьшит корреляцию $\rho$ между деревьями и предотвратит утечку целевой переменной.

\subsection*{2.}

При случайном выборе порога вместо оптимального каждое дерево строится хуже, так как меньше подстраивается под обучающую выборку.\\[1mm]
\textbf{Смещение} каждого отдельного дерева \textit{увеличивается}, поскольку порог выбирается не оптимально --- дерево хуже описывает закономерности в данных.\\[1mm]
\textbf{Разброс:} Случайность порогов делает деревья менее похожими друг на друга, то есть уменьшает корреляцию $\rho$:
\[
\text{Var}(a) = \rho\,\sigma^2 + \frac{1-\rho}{N}\sigma^2
\]
уменьшение $\rho$ при фиксированном $N$ снижает дисперсию ансамбля.\\[1mm]
Это типичный bias-variance trade-off. Рандомизация порогов немного увеличивает смещение каждого дерева, но значительно снижает разброс всего ансамбля за счёт декорреляции деревьев. Поскольку бэггинг не меняет смещение композиции, а снижает разброс, дополнительная рандомизация усиливает этот эффект.

\subsection*{3.}

Стандартная процедура градиентного бустинга где сдвиги вычисляются как антиградиент функции потерь:
\[
s_i = -\left.\frac{\partial L(y_i, z)}{\partial z}\right|_{z=a_{N-1}(x_i)}
\]
Для логистической функции потерь $L(y,z) = \log(1 + \exp(-yz))$ антиградиент равен:
\[
s_i = \frac{y_i}{1 + \exp(y_i a_{N-1}(x_i))}
\]
Значит, $|s_i| \leqslant 1$, при этом $|s_i|$ близок к 1, когда объект сильно ошибочно классифицирован, и близок к 0, когда объект уверенно верно классифицирован.\\[1mm]
\textbf{Модификация:} Вместо обучения на все сдвиги $s_i$, базовая модель обучается на $s_i [|s_i| \geqslant 0.1]$. Это означает, что если $|s_i| < 0.1$, то целевое значение для данного объекта обнуляется --- базовая модель не пытается его сдвинуть, иначе остается как есть.\\[1mm]

\textit{Первые базовые модели:}\\[1mm]
На старте композиция ещё не обучена ($a_0(x)$ --- простая константа), поэтому ошибки на большинстве объектов велики. Для логистических потерь:
\[
|s_i| = \frac{1}{1 + \exp(y_i a_0(x_i))}.
\]

При $a_0(x) = 0$ получаем $|s_i| = 0.5$ для всех объектов, что значительно больше порога 0.1. Следовательно, для первых моделей $[|s_i| \geqslant 0.1] = 1$ почти для всех объектов, и обучение практически не отличается от стандартного.\\

\textit{Поздние базовые модели:}\\[1mm]
По мере добавления деревьев композиция становится всё точнее. Для объектов, которые уже хорошо классифицированы (большой положительный отступ $y_i a_{N-1}(x_i) \gg 0$):
\[
|s_i| = \frac{1}{1 + \exp(y_i a_{N-1}(x_i))} \approx 0
\]

Для таких объектов $|s_i| < 0.1$ и их сдвиг обнуляется. Базовая модель перестаёт обращать на них внимание и фокусируется только на объектах, где $|s_i| \geqslant 0.1$ --- то есть на объектах, которые модель ещё не научилась классифицировать уверенно.

\textbf{Проблема:} На поздних итерациях объектов с $|s_i| \geqslant 0.1$ становится мало. Среди них окажутся преимущественно:
\begin{itemize}
\item Выбросы --- шумовые объекты с неверными метками
\item Объекты вблизи границы классов.
\end{itemize}

Базовая модель будет обучаться только на этих <<трудных>> объектах, что может привести к \textbf{переобучению на выбросах}: модель начнёт подстраиваться под шум, ухудшая качество на основной массе данных (по аналогии с проблемой экспоненциальной функции потерь)

\section*{Задача 3 (2.5 балла)}

Мы столкнулись с очень необычной задачей бинарной классификации: в выборке $\ell$ объектов, и лишь один из них относится к положительному классу. Модели обучать будет нелегко, но зато мы можем существенно упростить подсчёт метрик качества. Обозначим за $b(x)$ модель, выдающую уверенность в том, что $x$ относится к положительному классу.

Выполните следующие задания:

\begin{enumerate}
\item Пусть $k \in \{0, 1, \ldots, \ell-1\}$ --- позиция положительного объекта в выборке, ранжированной по уверенности модели. Иными словами, $k$ объектов имеют большее значение $b(x)$, чем положительный объект. Запишите формулу для AUC-ROC в зависимости от $k$.

\item Рассмотрим теперь новую метрику --- площадь под FNR/FPR кривой. По оси $x$ откладывается FNR $= \frac{FN}{FN+TP}$, по оси $y$ откладывается FPR $= \frac{FP}{FP+TN}$. Запишите формулу для зависимости площади под этой кривой от $k$.
\end{enumerate}

\subsection*{1.}

В выборке 1 положительный объект и $\ell - 1$ отрицательных. Всего пар (положительный, отрицательный):
\[
\ell_+ \cdot \ell_- = 1 \cdot (\ell - 1) = \ell - 1
\]
По аналитической формуле AUC-ROC:
\[
\text{AUC-ROC} = \frac{1}{\ell_+\,\ell_-}\sum_{i:\,y_i=1}\sum_{j:\,y_j=-1}\left([b(x_i) > b(x_j)] + \frac{1}{2}[b(x_i) = b(x_j)]\right)
\]
Поскольку $\ell_+ = 1$, внешняя сумма содержит лишь одно слагаемое. По условию, $k$ отрицательных объектов имеют оценку \textit{выше}, чем положительный объект. Значит, для них $b(x_i) < b(x_j)$ --- это неправильно упорядоченные пары. Правильно упорядоченных пар: $(\ell - 1) - k$.

Предполагая отсутствие равных оценок, что согласуется с условием:
\[
\text{AUC-ROC} = \frac{(\ell - 1) - k}{\ell - 1} = 1 - \frac{k}{\ell - 1}
\]

\textbf{Проверка:}
\begin{itemize}
\item При $k = 0$ --- положительный объект на первом месте: $\text{AUC} = 1$
\item При $k = \ell - 1$ --- положительный объект на последнем месте: $\text{AUC} = 0$
\end{itemize}

\subsection*{2.}

Кривая строится в осях: ось $x$ --- FNR, ось $y$ --- FPR. Их определения:
\[
\text{FNR} = \frac{FN}{FN+TP} = \frac{FN}{\ell_+}, \qquad
\text{FPR} = \frac{FP}{FP+TN} = \frac{FP}{\ell_-}
\]
Поскольку $\ell_+ = 1$ и $\ell_- = \ell - 1$:
\[
\text{FNR} \in \{0, 1\}, \qquad \text{FPR} \in \left\{0, \frac{1}{\ell-1}, \frac{2}{\ell-1}, \ldots, 1\right\}
\]
\textbf{Построение кривой.} Сортируем объекты по убыванию $b(x)$, при пороге выше максимального $b(x)$ все прогнозы отрицательные: $FP = 0$, $FN = 1$, тогда начальная точка: $(\text{FNR}, \text{FPR}) = (1, 0)$\\[1mm]
При понижении порога:
\begin{itemize}
\item Если проходим отрицательный объект: $FP$ растёт на 1, $\text{FPR}$ растёт на $\frac{1}{\ell-1}$ - шаг вверх
\item Если проходим положительный объект: $FN$ падает до 0, $\text{FNR}$ падает до 0 - шаг влево
\end{itemize}
В конце --- $(\text{FNR}, \text{FPR}) = (0, 1)$.
\begin{center}
\begin{tikzpicture}
\begin{axis}[
    width=10cm, height=8cm,
    xlabel={FNR}, ylabel={FPR},
    xmin=-0.05, xmax=1.1, ymin=-0.05, ymax=1.1,
    axis lines=middle,
    grid=both,
    xtick={0, 0.5, 1}, ytick={0, 0.4, 1},
    yticklabels={$0$, $\frac{k}{\ell-1}$, $1$}
]
\addplot[blue, very thick, ->] coordinates {(1, 0) (1, 0.4)};
\addplot[blue, very thick, ->] coordinates {(1, 0.4) (0, 0.4)};
\addplot[fill=blue!20, draw=none] coordinates {(0, 0) (1, 0) (1, 0.4) (0, 0.4) (0, 0)};
\node at (axis cs:0.5, 0.2) {\textbf{Площадь}};
\end{axis}
\end{tikzpicture}
\end{center}

\[
\text{Площадь под FNR/FPR кривой} = \frac{k}{\ell - 1} = 1 - \text{AUC-ROC}.
\]

\section*{Задача 4 (2.5 балла).}

Напомним, что разложение ошибки на шум, смещение и разброс выглядит так:
\[
L(\mu) = \underbrace{\mathbb{E}_{x,y}\!\left[\bigl(y - \mathbb{E}[y\mid x]\bigr)^2\right]}_{\text{шум}} +
\underbrace{\mathbb{E}_x\!\left[\bigl(\mathbb{E}_X[\mu(X)] - \mathbb{E}[y\mid x]\bigr)^2\right]}_{\text{смещение}} +
\underbrace{\mathbb{E}_x\!\left[\mathbb{E}_X\!\left[\bigl(\mu(X) - \mathbb{E}_X[\mu(X)]\bigr)^2\right]\right]}_{\text{разброс}}
\]

Допустим, у нас есть обучающая выборка из двух объектов, в которой данные описываются одним бинарным признаком $x$, который принимает значение 1 с вероятностью $p$ и значение 0 с вероятностью $1-p$. Целевая переменная имеет нормальное распределение со средним $x$ и дисперсией $\sigma^2$, $y \sim \mathcal{N}(x, \sigma^2)$. Метод обучения выдаёт модель, предсказывающую среднее по выборке: $\mu(X)(x) = \frac{x_1 + x_2}{2}$. Найдите шум, смещение и разброс описанного метода обучения.

\subsection*{Решение}


\subsubsection*{Истинная зависимость $\mathbb{E}[y\mid x]$}

По условию, $y \sim \mathcal{N}(x, \sigma^2)$, следовательно $\mathbb{E}[y \mid x] = x$\\

\subsubsection*{Шум $\text{Var}(y\mid x)$}

\[
\text{Шум} = \mathbb{E}_{x,y}\!\left[(y - \mathbb{E}[y\mid x])^2\right] = \mathbb{E}_x\!\left[\text{Var}(y\mid x)\right]
\]
Из условия $y \sim \mathcal{N}(x, \sigma^2)$:
\[
\text{Var}(y\mid x) = \sigma^2
\]
Эта величина не зависит от $x$, поэтому:
\[
\text{Шум} = \sigma^2
\]

\subsubsection*{Средний прогноз модели $\mathbb{E}_X[\mu(X)(x)]$}
\[
\mu(X)(x) = \frac{x_1 + x_2}{2}
\]
Прогноз модели \textbf{не зависит} от тестового объекта $x$, ведь он определяется только обучающей выборкой $(x_1, x_2)$\\[1mm]
Берём матожидание по случайности обучения ($x_1$ и $x_2$ --- независимые случайные величины с распределением $\text{Ber}(p)$):
\[
\mathbb{E}_X[\mu(X)(x)] = \mathbb{E}_{x_1, x_2}\!\left[\frac{x_1 + x_2}{2}\right] = \frac{\mathbb{E}[x_1] + \mathbb{E}[x_2]}{2}
\]
Поскольку $x \sim \text{Ber}(p)$, то $\mathbb{E}[x] = p$:
\[
\mathbb{E}_X[\mu(X)(x)] = \frac{p + p}{2} = p
\]
\subsubsection*{Смещение}
\[
\text{Bias}(x) = \mathbb{E}_X[\mu(X)(x)] - \mathbb{E}[y\mid x] = p - x
\]
Квадрат смещения:
\[
\text{Bias}^2(x) = (p - x)^2
\]
Усредняем по тестовому $x \sim \text{Ber}(p)$:
\[
\mathbb{E}_x[\text{Bias}^2(x)] = \mathbb{E}_x[(x - p)^2] = \text{Var}(x)
\]
Для бернуллиевской случайной величины $\text{Var}(x) = p(1-p)$:
\[
\text{Смещение}^2 = p(1-p)
\]

\subsubsection*{Разброс $\text{Var}_X(\mu(X)(x))$}

\[
\text{Var}_X(\mu(X)(x)) = \text{Var}_{x_1, x_2}\!\left(\frac{x_1 + x_2}{2}\right) = \frac{1}{4}\bigl(\text{Var}(x_1) + \text{Var}(x_2)\bigr)
\]
Поскольку $x_1, x_2$ независимы и одинаково распределены, $\text{Var}(x_i) = p(1-p)$:
\[
\text{Var}_X(\mu(X)(x)) = \frac{1}{4}\bigl(p(1-p) + p(1-p)\bigr) = \frac{p(1-p)}{2}
\]
Разброс не зависит от $x$, поэтому усреднение по $x$ не меняет результат:
\[
\text{Разброс} = \frac{p(1-p)}{2}
\]

\[
\text{Шум} = \sigma^2, \qquad
\text{Смещение}^2 = p(1-p), \qquad
\text{Разброс} = \frac{p(1-p)}{2}
\]
\begin{itemize}
\item \textbf{Шум $= \sigma^2$:} неустранимая ошибка, обусловленная случайностью $y$ при фиксированном $x$

\item \textbf{Смещение $= p(1-p)$:} Средний прогноз модели ($= p$) не совпадает с истинной зависимостью ($= x$). Это объясняется тем, что модель предсказывает среднее по \textit{обучающей выборке}, которое не зависит от конкретного тестового объекта $x$. Смещение тем больше, чем более <<сбалансирован>> признак ($p$ близко к $0.5$)

\item \textbf{Разброс $= p(1-p)/2$:} Прогноз модели зависит от обучающей выборки через $x_1 + x_2$. При $N$ объектах в выборке разброс был бы $p(1-p)/N$

\item Полная ошибка:
\[
L(\mu) = \sigma^2 + p(1-p) + \frac{p(1-p)}{2} = \sigma^2 + \frac{3}{2}p(1-p)
\]
\end{itemize}

\end{document}